\section{后端切换：\texttt{kconfig.h}}

三种后端已全部实现，通过 \texttt{kconfig.h} 的编译期开关选择：

\codefile{src/include/kernel/kconfig.h}
\begin{lstlisting}[style=cstyle]
/*  0 = sorted-array  — 排序数组，简单稳定，早期默认
 *  1 = rbtree         — 红黑树，O(log n) 全操作
 *  2 = maple-tree     — B+ tree 变体，高扇出缓存友好
 */
#define KCONFIG_VMA_BACKEND  0
\end{lstlisting}

\texttt{kmain.c} 中根据宏值选择后端：

\codefile{src/kernel/core/kmain.c}
\begin{lstlisting}[style=cstyle]
#if KCONFIG_VMA_BACKEND == 0
    vma_register_backend(vma_sorted_array_get_ops());
#elif KCONFIG_VMA_BACKEND == 1
    vma_register_backend(vma_rbtree_get_ops());
#elif KCONFIG_VMA_BACKEND == 2
    vma_register_backend(vma_maple_get_ops());
#endif
\end{lstlisting}

修改 \texttt{KCONFIG\_VMA\_BACKEND} 后 \texttt{make iso}，
用 \texttt{vma count} 命令确认当前加载的后端名称。
\texttt{vma test} 自测试对所有三种后端均通过。

% ============================================================================

